\documentclass[11pt,a4paper]{article}

\usepackage[utf8]{inputenc}

\usepackage[T1]{fontenc}

\usepackage[polish]{babel}

\usepackage{lmodern}

\usepackage{geometry}

\geometry{a4paper,margin=27mm}

\usepackage{microtype}

\usepackage[hidelinks]{hyperref}

% Portfolio: MSC-AR-P14; source: prospective ARC-XR research scope.

\title{Empatyczny awatar AI adaptujący komunikację do pobudzenia użytkownika}

\author{Mateusz Pomianek}

\date{Wersja 0.3.0, 21 listopada 2026}

\begin{document}

\maketitle

\begin{abstract}

Przedmiotem opracowania jest Empatyczny awatar AI adaptujący komunikację do pobudzenia użytkownika. Czy awatar dostosowujący styl komunikacji do pobudzenia poprawia rozumienie wskazówek w XR? Projekt obejmuje porównanie stałego i adaptowanego stylu komunikatów awatara AI przy tej samej treści oraz liczbie wskazówek. Głównym punktem końcowym jest poprawność odczytania wskazówki, ocena odbioru i liczba błędnych adaptacji. Opis ma charakter wstępny i nie przedstawia niezweryfikowanych wyników.

\end{abstract}

\section{Wprowadzenie}

Rozwój adaptacyjnych systemów XR stwarza możliwość badania interakcji człowieka z informacją przestrzenną, ale samo zastosowanie AR, VR lub modeli sztucznej inteligencji nie przesądza o przewadze poznawczej czy ergonomicznej. W rozpatrywanym zagadnieniu, Empatyczny awatar AI adaptujący komunikację do pobudzenia użytkownika, kluczowe jest rozdzielenie właściwości systemu od zmienności zachowania użytkownika. Porównanie należy zaprojektować tak, aby mierzyć jednoznaczne zjawisko, a nie tylko ogólną ocenę technologii.

\section{Cel pracy i pytanie badawcze}

Czy awatar dostosowujący styl komunikacji do pobudzenia poprawia rozumienie wskazówek w XR? Projekt wymaga krytycznego przeglądu aktualnego stanu badań oraz zdefiniowania nowości w stosunku do innych tematów portfolio XR.

\section{Zakres i odrębny wkład naukowy}
Adaptacja sposobu wypowiedzi awatara XR do sygnału pobudzenia; zawartość instrukcji i liczba komunikatów pozostają identyczne. Pobudzenie nie jest etykietą emocji. Przyczynowy wpływ formy komunikacji przy równej treści, nie jakość wykrywania emocji.
\section{Hipotezy badawcze}
\paragraph{H1.} Adaptowany styl komunikatu awatara zwiększa poprawność zrozumienia wskazówki w porównaniu ze stylem stałym.
\paragraph{H2.} Zysk nie wynika wyłącznie z wyższej oceny empatii, gdyż utrzymuje się po kontroli treści, szybkości i częstotliwości komunikatów.
\paragraph{H3.} Przy niepewnym sygnale pobudzenia adaptacja może pogarszać skuteczność, więc jawna abstencja zmniejsza niekorzystne błędy dopasowania.
\section{Projekt badania i plan eksperymentów}
\subsection{A. Operacjonalizacja i kwalifikacja danych}
Adaptacja sposobu wypowiedzi awatara XR do sygnału pobudzenia; zawartość instrukcji i liczba komunikatów pozostają identyczne. Pobudzenie nie jest etykietą emocji. Jednostka wnioskowania: uczestnik × kompletna instrukcja lub epizod dialogu. W protokole zamrozić definicję i mianownik wyniku oraz źródło prawdy odniesienia, zgodność znaczników czasu, prawa do danych i zasady rekrutacji.
\subsection{B. Eksperyment porównawczy}
Porównać warunek stały, adaptację ze zweryfikowanym sygnałem i kontrolę o tej samej długości wypowiedzi; nie dopuszczać generowania innych instrukcji. Estymand pierwszorzędowy: różnica prawdopodobieństwa poprawnego wykonania wskazówki dla adaptowanego vs stałego stylu komunikatu. Przed testem zdefiniować randomizację, próbę, punkt końcowy i warunki wykluczenia bez selekcji wyniku post hoc.
\subsection{C. Kontrole i analiza odporności}
Oddzielnie zebrać deklarowaną ocenę komunikatu i obiektywny wynik zadania oraz niezależnie kontrolować trafność sygnału. Analizować błędy w zależności od braku danych i odsetka odrzuconych adaptacji. W rejestrze uwzględnić braki, wstrzymania, opóźnienia i niepełne próby zamiast usuwać je z analizy.
\subsection{D. Walidacja na niezależnym materiale}
Weryfikować na niezależnych sesjach lub osobach, bez dostępu modelu do danych końcowych. Model Bernoulliego wykonania wskazówki z uczestnikiem i treścią instrukcji; analiza warunkowa przy opóźnieniu. Dokumentować wyniki niesprzyjające hipotezie oraz rzeczywiste granice przeniesienia.
\section{Estymand, analiza statystyczna i decyzje}
różnica prawdopodobieństwa poprawnego wykonania wskazówki dla adaptowanego vs stałego stylu komunikatu. Jednostka: uczestnik × kompletna instrukcja lub epizod dialogu. Model Bernoulliego wykonania wskazówki z uczestnikiem i treścią instrukcji; analiza warunkowa przy opóźnieniu. H1 jest kandydatem hipotezy potwierdzającej, a H2/H3 pozostają eksploracyjne do formalnego planu wielokrotnych porównań. Liczebność uczestników, margines, próg i moc lub precyzja wymagają pilotażu oraz zamrożenia przed testem końcowym.
\section{Stan badań i granica nowości}
\begin{itemize}
\item Proactive LMM-based AR assistance exists; evaluate the behavioural difference under a fixed number of interventions. \cite{xrprior1}.
\item Human reliance on AR-provided AI instructions is already studied; measuring trust alone does not establish correct guidance. \cite{xrprior2}.
\end{itemize}
Przyczynowy wpływ formy komunikacji przy równej treści, nie jakość wykrywania emocji. Odrębność od wcześniejszych publikacji i innych tematów programu należy sprawdzić przed rejestracją badania.
\section{Warunki bezpieczeństwa, sprzętu i zgodności}
Jest to plan prospektywny, nie dowód wykonania eksperymentu lub zatwierdzonej etyki. Udział osób, rejestrowanie biometrii i interwencje wymagają uprzedniej procedury etycznej i świadomej zgody. Nie wnioskować o emocjach ani zdrowiu z samych sygnałów. Standardowy Quest 3 nie ma okulometrii HMD, a Tobii Eye Tracker 5 służy do stanowiska ekranowego; każde połączenie zweryfikować przed protokołem.
\section{Materiał i metody}

Projekt obejmuje porównanie stałego i adaptowanego stylu komunikatów awatara AI przy tej samej treści oraz liczbie wskazówek. Wymagane jest zdefiniowanie jednostki analizy, zasad doboru przypadków, wersji scenariusza i źródeł pomiaru. Rozkład warunków oraz reguły wykluczenia prób należy określić przed pomiarem. W badaniach z udziałem ludzi niezbędna jest uprzednia właściwa ocena etyczna i świadoma zgoda. W testach wyłącznie symulacyjnych należy oddzielić prawdę symulacyjną od obserwacji empirycznej.

\subsection{Wynik główny i kryteria jakości}

Poprawność odczytania wskazówki, ocena odbioru i liczba błędnych adaptacji. Pomiar powinien zachowywać jawne jednostki, znaczniki czasu i reguły oceny poprawności. Przy modelach predykcyjnych stosować rozłączne partycje oceny; przy eksperymentach interwencyjnych oddzielić skuteczność predykcji od przyczynowego wpływu interfejsu.

\section{Plan analizy}

Analiza efektu interfejsu oddzielona od błędu modelu stanu afektywnego, z kontrolą saliencji przekazu. W wynikach należy raportować wielkość efektu lub błąd estymacji z przedziałem niepewności, liczbą obserwacji i odsetkiem odrzuconych prób. Analizy dodatkowe wymagają jasnego oznaczenia jako eksploracyjne. Nie wybierać progów decyzyjnych ani modelu na podstawie obserwacji ze zbioru testowego.

\section{Wyniki}

W obecnej wersji nie raportuje się wyników pomiarów, tabel statystycznych ani parametrów wyznaczonych z nieudokumentowanych danych. To kierunek badawczy, a nie zatwierdzony protokół badania ani zgłoszony artykuł.

\section{Dyskusja}

Interpretacja przyszłych wyników musi uwzględnić alternatywne mechanizmy i zależność wniosków od jakości danych, warunków zadania oraz ograniczeń sprzętu. Wartość potencjalnego rozwiązania zależy nie tylko od trafności modelu, lecz także od jego opóźnienia, stabilności i wpływu na zachowanie człowieka. Wyniki testów syntetycznych nie mogą być automatycznie uogólniane na populację uczestników.

\section{Ograniczenia badania}

Błędna estymacja afektu może pogarszać interakcję, a antropomorfizacja zmieniać deklarowane zaufanie. Granice nowości i trafności ekologicznej wymagają niezależnej weryfikacji. Żadna część szkicu nie dowodzi przeprowadzenia zbierania danych, uzyskania zgód etycznych ani gotowości manuskryptu do zgłoszenia.

\section{Wnioski}

Przedstawiono problem, kontrast i sposób jego planowanej oceny. Rozstrzygnięcie postawionego pytania wymaga przygotowania protokołu, zweryfikowanych źródeł oraz poprawnie udokumentowanych wyników. Bez tych materiałów nie można sformułować konkluzji empirycznych.

\begin{thebibliography}{9}
\bibitem{xrprior1} Sensible Agent: A Framework for Unobtrusive Interaction with Proactive AR Agents (2025), UIST 2025. \url{https://discovery.ucl.ac.uk/id/eprint/10224695/}.
\bibitem{xrprior2} An Investigation of Trust and Overreliance on Conversational Instructions in Augmented Reality (2026), CUI 2026. \url{https://doi.org/10.1145/3816046.3816296}.
\end{thebibliography}
\end{document}

